\section{Open-Source LLMs and Local Deployment}
The democratization of large language models has been accelerated by the rise of open-source initiatives, allowing researchers and organizations to deploy high-performance models on their own infrastructure. Local deployment is particularly critical for industrial MLOps, where data privacy, latency, and operational cost are paramount concerns.

\subsection{The Rise of Open-Source Model Families}
Models such as Meta’s \textbf{Llama 3} and \textbf{Mistral} have demonstrated that smaller, well-trained models can approach the performance of massive proprietary APIs on many reasoning and coding tasks \cite{mistral_7b, touvron2023llama}. These models are often distributed in quantized formats (e.g., GGUF), which compress the model weights from 16-bit floating point to 4-bit or 8-bit integers. This reduction in memory footprint allows state-of-the-art models (even those with 8B to 70B parameters) to run on consumer-grade hardware or local workstations without specialized GPU clusters.

\subsection{Ollama: Streamlining Local Inference}
\textbf{Ollama} is a lightweight, open-source framework designed to simplify the local deployment and management of LLMs \cite{ollama_github_repo}. It packages model weights, configurations, and inference engines into a single, easy-to-use interface. By utilizing the \texttt{llama.cpp} backend, Ollama provides high-performance inference through hardware acceleration (Metal on macOS, CUDA on NVIDIA, and ROCm on AMD).

Key features of Ollama integrated into this work include:
\begin{enumerate}
    \item \textbf{Model Versioning:} Simple commands (e.g., \texttt{ollama pull mistral}) allow for quick swapping between different model variants during testing.
    \item \textbf{Custom Modelfiles:} Developers can define specialized system prompts and parameters (e.g., temperature, context window) directly in a declarative Modelfile.
    \item \textbf{REST API Integration:} Ollama exposes a standard local API, allowing LangChain and LangGraph to consume local models as if they were remote cloud endpoints, facilitating a "hybrid" architecture where sensitive code generation happens locally while non-sensitive research tasks can use larger cloud models.
\end{enumerate}
In local deployment scenarios, tools like Ollama serve as the primary inference backend, ensuring that all model "surgery" and firmware configuration logic remains within the user's local network, thereby preserving intellectual property and reducing external dependencies.
